\section{Discussion}
\label{sec:discussion}

The current implementation prioritizes a reproducible and modular architecture over immediate benchmark-level performance. This choice is appropriate for an initial research scaffold, but several limitations remain. First, the FEMNIST ingestion pipeline is simplified through feature compression, so the model does not yet exploit the full spatial structure of the original images \cite{leaf}. Second, the quantum model is intentionally compact and may not capture the expressivity required for high-accuracy classification. Third, the unlearning metric based on QFI should be interpreted as a geometric indicator rather than a complete proof of erasure \cite{qfi,unlearning_survey_2024}.

Despite these limitations, the project establishes a clear experimental path: it separates training and unlearning, integrates GPU-aware execution in PennyLane, and introduces a privacy-oriented evaluation protocol based on forget-set performance, retain-set performance, membership inference risk, and QFI sensitivity. From the perspective of research contribution, the main value of the current repository lies in providing a clean software scaffold that can be extended toward stronger quantum ansätze, richer federated aggregation strategies, and more formal unlearning guarantees.

A further open problem is the lack of a full FEMNIST ingestion pipeline directly from LEAF. Until that is implemented, the current `npz` intermediary should be treated as an experimental convenience rather than as a final benchmark protocol. Future work should also validate whether the QFI-based signal correlates with stronger empirical unlearning criteria under multiple random seeds and client-dropout scenarios \cite{federated_unlearning_active_forgetting,forgettable_federated_linear_learning,qfl_foundations_2023}.
